% !TEX root = ../linnik399.tex
\section{The case tree}\label{sec:casetree}

This section describes the finite case analysis of the middle range $0.1\le\lambda_1<1.5$ and
proves that it covers every configuration of zeros (\cref{thm:cover}). The analytic facts it
uses are the zero-location results of \cref{sec:location}, the costs of \cref{sec:costs}, the far
budget of \cref{sec:far} and the near rows of \cref{sec:rows}. There are two tasks:
showing that the cases cover every possible configuration, and showing that a
configuration in a case gives a feasible point of its program. These are proved
separately before being combined in \cref{thm:cover}. \Cref{tab:census}, after \cref{subsec:refinement}, lists the levels of the
analysis, with their numbers and the places where they are defined and checked.

\subsection{Configurations}\label{subsec:configurations}
\begin{definition}[Zero configuration]\label{def:configuration}
Fix a sufficiently large $q$. Its \emph{zero configuration} is the multiset $\mathcal Z$ of pairs $(\chi,\rho)$ with
$\chi\ne\chi_0$ a character modulo $q$, $\rho\in R(l)$ and $L(\rho,\chi)=0$, each pair repeated
according to the multiplicity of $\rho$.
\end{definition}

Minima over empty sets are $+\infty$. In the middle range the first zero exists.
The quantities of \cref{sec:notation} are read off from $\mathcal Z$: the first zero $(\chi_1,\rho_1)$,
the first family $\mathcal F_1$ with $n=|\mathcal F_1|$, the parameters $\lambda_2,\lambda_3$ and the family
$\mathcal F_2$ (\cref{def:families}), and the distinguished occurrences, $\lambda'$ and an admissible
$\rho'$ (\cref{def:additionalzero}); for type $\mathrm{rc}$, $\mu_1>0$. In addition we use the
following.
\begin{itemize}
  \item The \emph{height-one parameter} of $\chi$ is $\nu(\chi)=\min\{\lambda_\rho:(\chi,\rho)\in\mathcal
    Z,\ |\gamma_\rho|\le1\}$, and the \emph{$T^*$-parameter} is $\nu_{T^*}(\chi)$, defined in the same
    way with $|\gamma_\rho|\le T^*$ (\cref{def:representative}). Since $L(\rho,\chi)=0$ if and only if
    $L(\overline\rho,\bchi)=0$, both are functions of the family. We put
    $\nu_*=\min\{\nu(\mathcal F):\mathcal F\ne\mathcal F_1\}$.
  \item The configuration is \emph{inside} if $\rho_1\in R_B$ and \emph{outside} otherwise. Since
    $\lambda_1<1.5<C_B$ in the middle range, outside means $|\mu_1|>C_B$. Type $\mathrm{rr}$ is
    always inside.
\end{itemize}
The definitions give $\lambda_1\le\lambda'$, $\lambda_1\le\lambda_2\le\lambda_3$, and
$\nu_{T^*}(\chi)\ge\nu(\chi)\ge\lambda_2$ for $\chi\notin \mathcal F_1$.

\subsection{Specifications}\label{subsec:specifications}

A specification records interval information about a configuration. It may describe
no actual configuration; in that event the corresponding case can be excluded.
Its two lower bounds $l_2$ and $r$ have different scopes: $l_2$ applies throughout
$R(l)$, whereas $r$ applies only to the height-one representatives.

\begin{definition}[Specification]\label{def:specification}
A \emph{specification} is a tuple
\[
  \mathfrak s=(\mathfrak t,[a,b],p,\mathrm{gap},l_2,r,\mathrm{res},\mathrm{loc})
\]
with the following data:
\begin{itemize}
  \item a type $\mathfrak t\in\{\mathrm{rr},\mathrm{rc},\mathrm{complex}\}$ and a
    first-zero cell $[a,b]$, where $a<b$;
  \item a lower bound $p$ for $\lambda'$, and optionally an interval
    $\mathrm{gap}=[\lo',\hi']$ with $p\le\lo'<\hi'\le\infty$, called a \emph{gap}
    (used only for type $\mathrm{complex}$);
  \item a global second-family lower bound $l_2$ and a height-one lower bound $r$,
    with $l_2\le r\le3$;
  \item either no reservation, or a pair $\mathrm{res}=(\hi_2,n_2)$ with
    $r<\hi_2\le2$ and $n_2\in\{1,2\}$; in the latter case put $\lo_2=r$;
  \item a height class $\mathrm{loc}\in\{\mathrm{inside},\mathrm{outside}\}$.
\end{itemize}
The \emph{configuration set} $C(\mathfrak s)$ consists of the configurations for which
some admissible choices of the first zero and the minimizing families satisfy all of
the following:
\begin{enumerate}[(C1)]
  \item the type is $\mathfrak t$ and $\lambda_1\in[a,b]$;
  \item $\lambda'\ge p$, and $\lambda'\in[\lo',\hi']$ if a gap is present (with $\lambda'=\infty$ allowed
    when $\hi'=\infty$);
  \item $\lambda_2\ge l_2$;
  \item if $\mathrm{res}$ is none, then $\nu(\mathcal F)\ge r$ for every family $\mathcal F\ne\mathcal F_1$; if
    $\mathrm{res}=(\hi_2,n_2)$, then $\nu_*\in[r,\hi_2]$ and some family $\mathcal F_{\mathrm{res}}\ne\mathcal F_1$ with
    $\nu(\mathcal F_{\mathrm{res}})=\nu_*$ has exactly $n_2$ characters (so it is a real character if $n_2=1$ and
    a pair of conjugate nonreal characters if $n_2=2$);
  \item the configuration is inside if $\mathrm{loc}=\mathrm{inside}$ and outside if $\mathrm{loc}=\mathrm{outside}$.
\end{enumerate}
In the reserved case, (C4) implies $\nu(\mathcal F)\ge r$ for every $\mathcal F\ne\mathcal F_1$. We call $\mathcal F_{\mathrm{res}}$ the
\emph{reserved family} and its characters \emph{reserved}; the characters outside $\mathcal F_1$ and
$\mathcal F_{\mathrm{res}}$ are \emph{ordinary}.
\end{definition}

In a formula containing $\lo'$, that argument is omitted if no gap is present.
Likewise, reserved-family terms are absent when there is no reservation.

The reserved family is a family with the least height-one parameter among the families other than
the first; it is not assumed to be the family $\mathcal F_2$ that attains $\lambda_2$, which may have its
nearest zero at a height above $1$. The following three facts are used in every leaf.

\begin{lemma}[Consequences of reserving a family]\label{lem:familyfacts}
Let $\mathfrak s$ be a specification and fix a configuration and admissible choices
witnessing its membership in $C(\mathfrak s)$.
\begin{enumerate}[(i)]
  \item \emph{(Count.)} At most two characters $\chi\notin \mathcal F_1$ have $\nu_{T^*}(\chi)<\lambda_3$.
  \item \emph{(Drop.)} If $\mathfrak s$ is reserved, then $\nu_{T^*}(\chi)\ge\nu(\chi)\ge\lambda_3$ for every
    ordinary character $\chi$.
  \item \emph{(Cap.)} If $\mathfrak s$ is reserved, then $\lambda_2\le\nu_*\le\hi_2$.
\end{enumerate}
\end{lemma}

\begin{proof}
(i) Such a character has a zero in $R(l)$ with parameter below $\lambda_3$, so it lies in $\mathcal F_2$, and
$|\mathcal F_2|\le2$. (ii) Let $\mathcal F$ be the family of an ordinary character. If $\mathcal F\ne\mathcal F_2$, every zero of $\mathcal F$
in $R(l)$ has parameter at least $\lambda_3$. If $\mathcal F=\mathcal F_2$, then $\mathcal F_{\mathrm{res}}\notin\{\mathcal F_1,\mathcal F_2\}$, so
$\nu(\mathcal F_{\mathrm{res}})\ge\lambda_3$, and $\nu(\mathcal F)\ge\nu_*=\nu(\mathcal F_{\mathrm{res}})$ by minimality. (iii) The
height-one representative of $\mathcal F_{\mathrm{res}}$ is a zero in $R(l)$ of a character outside $\mathcal F_1$.
\end{proof}

\subsection{The source cover}\label{subsec:sourcecover}
The roots of the case tree are $2768$ specifications, each used with the height class
$\mathrm{inside}$, and the $1685$ of them whose type is not $\mathrm{rr}$ used also with the height
class $\mathrm{outside}$: $4453$ roots in all. They are constructed from the zero-location results of
\cref{sec:location} in four steps.
\begin{enumerate}[(1)]
  \item \emph{Parent rows.} There are $58$ \emph{parent rows} $(\mathfrak t,A_j,B_j,p_j,r_j)$: $33$ of type
    $\mathrm{rr}$ tiling $[0.1,1.5]$, $13$ of type $\mathrm{rc}$ tiling $[0.628,1.5]$, and $12$ of type
    $\mathrm{complex}$ tiling $[0.44,1.5]$. Each is the statement
    \begin{equation}\tag{T2}\label{eq:T2}
      \text{type }\mathfrak t,\ \lambda_1\in[A_j,B_j]\ \Longrightarrow\ \lambda'\ge p_j\ \text{and}\ \lambda_2\ge r_j,
    \end{equation}
    which is proved in \cref{prop:parentrows} from the tables of Heath-Brown and Xylouris.
  \item \emph{Base cells.} Each parent interval is cut into cells of width $0.01$, giving $340$
    \emph{base cells}.
  \item \emph{Cells and gaps.} The base cells are tiled by $478$ first-zero cells in all. A cell
    $[\lo,\hi]$ of the parent $j$ receives numbers $p\le\max(p_j,\lo)$ and $l_2\le\max(r_j,\lo)$; this is
    valid since $\lambda',\lambda_2\ge\lambda_1\ge\lo$. A cell of type
    $\mathrm{complex}$ may be \emph{partitioned} by gaps $[p,m_1],[m_1,m_2],\dots,[m_k,\infty]$ for
    $\lambda'$. The cells that are not partitioned, and the gaps of those that are, are the $684$
    \emph{gap cases}.
  \item \emph{Records.} Each gap case is either \emph{unreserved}, with one specification with
    $r=l_2$ and no reservation, or \emph{reserved} with a chain $l_2=u_0<u_1<\dots<u_m=e\le2$,
    with specifications $(\mathrm{res}=(u_{j+1},n_2),\,r=u_j)$ for $n_2=1,2$ and $0\le j<m$,
    followed by an unreserved \emph{tail} specification with $r=e$.
\end{enumerate}
By type there are $1083$ specifications of type $\mathrm{rr}$ ($195$ unreserved, $444$ reserved
with $n_2=1$ and $444$ with $n_2=2$), $115$ of type $\mathrm{rc}$ (all unreserved and without gap)
and $1570$ of type $\mathrm{complex}$.

\begin{theorem}[Coverage by the root cases]\label{thm:sourcecover}
Every configuration with $0.1\le\lambda_1<1.5$ lies in $C(\mathfrak s)$ for at least one of the $4453$
root specifications $\mathfrak s$.
\end{theorem}

\begin{proof}
Let $\mathfrak t$ be the type. By \cref{prop:firstlower}, $\lambda_1>0.628$ if $\mathfrak t=\mathrm{rc}$ and
$\lambda_1>0.44$ if $\mathfrak t=\mathrm{complex}$, so $\lambda_1$ lies in some cell of type $\mathfrak t$ (the
cells are closed, and a boundary value lies in two cells). By \eqref{eq:T2} for the parent of the
cell, $\lambda'\ge p$ and $\lambda_2\ge l_2$. If the cell is partitioned, some gap contains $\lambda'$.
In the corresponding gap case, if it is unreserved, every family $\mathcal F\ne\mathcal F_1$ has $\nu(\mathcal F)\ge\lambda_2\ge
l_2=r$. If it is reserved with chain $u_0<\dots<u_m$, then $\nu_*\ge\lambda_2\ge u_0$; if $\nu_*<e$
we choose $j$ with $\nu_*\in[u_j,u_{j+1}]$ and a minimizing family, which is a real character or a
nonreal pair, and the configuration lies in the corresponding record; otherwise it lies in the
tail record. Finally the height class is determined by $|\mu_1|$.
\end{proof}

\subsection{Refinement trees}\label{subsec:refinement}
Each root is the root of a finite \emph{refinement tree}. Every internal node carries a
specification (computed from the root, never read from the data) and is of one of the following
kinds. In each case every configuration of the node lies in the configuration set of one of its
children, or the node has no configuration at all.
\begin{itemize}
  \item \emph{First-zero split} at $a<m<b$: children with cells $[a,m]$ and $[m,b]$, all other data
    kept. Exhaustive, since $\lambda_1\le m$ or $\lambda_1\ge m$.
  \item \emph{Second-family split} at $r<m<\hi_2$ (reserved nodes): children $(\hi_2:=m)$ and
    $(r:=m,\ \lo_2:=m)$. Exhaustive, since $\nu_*\le m$ or $\nu_*\ge m$, and in the second case every
    family $\mathcal F\ne\mathcal F_1$ has $\nu(\mathcal F)\ge\nu_*\ge m$.
  \item \emph{Gap split} at $\lo'<m<\hi'$ (type $\mathrm{complex}$): children with gaps $[\lo',m]$ and
    $[m,\hi']$. Exhaustive.
  \item \emph{Location update}: a lower bound $\lambda_2>h$ valid on the whole cell
    (\cref{prop:realrows,prop:complexrows}; \cref{lem:updates}) replaces $l_2$ and $r$ by $\max(l_2,h)$ and
    $\max(r,h)$, and excludes the node if it is reserved with $\hi_2\le h$. Sound, since every family
    $\mathcal F\ne\mathcal F_1$ has $\nu(\mathcal F)\ge\lambda_2$.
  \item \emph{Exclusion} of a reserved interval or a gap by a lower bound for $\lambda_2$ or
    $\lambda'$ (\cref{prop:polyrows,prop:realfirstsecond}), or of a whole node by the positivity
    argument of \cref{prop:positivity}.
\end{itemize}
Before a leaf model is built, all applicable location updates are applied once more to its
specification; this only adds valid implications. None of these nodes depends on $L$. The terminal
nodes are the \emph{leaves}. The inside trees have $3146$ leaves. Of these, $233$ (in $199$ roots) are
stored in a separate file and are called \emph{supplementary leaves}; the reason is historical
(\cref{rem:history}). Every leaf, supplementary or not, receives the same model. The
outside trees have $1642$ leaves. In all there are $4788$ leaves; $343$ inside roots and $70$
outside roots have none, because all their configurations are excluded.

\begin{table}[ht]
\centering
\small
\begin{tabular}{lrr}
\toprule
Node & Inside & Outside\\
\midrule
first-zero / second-family / gap splits & 1 / 408 / 593 & 0 / 27 / 0\\
location updates, real (restrict / exclude) & 225 / 75 & --\\
location updates, complex (restrict / exclude) & 11 / 0 & --\\
exclusions by lower bounds for $\lambda_2$ or $\lambda'$ & 352 & 70\\
exclusions by positivity & 197 & 0\\
leaves & 3146 & 1642\\
\bottomrule
\end{tabular}
\caption{Nodes of the refinement trees.}
\label{tab:nodes}
\end{table}

\begin{table}[tp]
\centering
\footnotesize
\setlength{\tabcolsep}{4pt}
\begin{tabular}{@{}>{\raggedright\arraybackslash}p{\dimexpr0.145\linewidth-6.4pt\relax}
                >{\raggedleft\arraybackslash}p{\dimexpr0.115\linewidth-6.4pt\relax}
                >{\raggedright\arraybackslash}p{\dimexpr0.37\linewidth-6.4pt\relax}
                >{\raggedright\arraybackslash}p{\dimexpr0.165\linewidth-6.4pt\relax}
                >{\raggedright\arraybackslash}p{\dimexpr0.205\linewidth-6.4pt\relax}@{}}
\toprule
Level & Number & What it partitions or represents & Where defined & How it is checked\\
\midrule
Parent rows & $58$ &
  Published implications: type $\mathfrak t$ and $\lambda_1\in[A_j,B_j]$ give $\lambda'\ge p_j$,
  $\lambda_2\ge r_j$ ($33$ $\mathrm{rr}$, $13$ $\mathrm{rc}$, $12$ $\mathrm{complex}$) &
  \cref{prop:parentrows} &
  Printed tables; Heath-Brown's Tables~4 and~7 recomputed\\
Base cells & $340$ &
  Parent intervals cut into cells of width $0.01$ ($145$, $89$, $106$) &
  \cref{subsec:sourcecover} &
  Validated in the replay\\
First-zero cells & $478$ &
  Tiling of the base cells ($195$, $115$, $168$), with bounds $p\le\lambda'$ and $l_2\le\lambda_2$ &
  \cref{subsec:sourcecover} &
  Validated in the replay\\
Gap cases & $684$ &
  $422$ cells without a gap; $262$ gaps for $\lambda'$ partitioning $56$ $\mathrm{complex}$ cells &
  \cref{subsec:sourcecover} &
  Validated in the replay\\
Specifications & $2768$ &
  Gap cases with a second-family record: $1083$ $\mathrm{rr}$ ($195$ unreserved, $444+444$
  reserved), $115$ $\mathrm{rc}$, $1570$ $\mathrm{complex}$ &
  \cref{def:specification} &
  Validated in the replay\\
Root cases & $4453$ &
  $2768$ specifications inside, and the $1685$ not of type $\mathrm{rr}$ also outside &
  \cref{thm:sourcecover} &
  Hand proof; replay\\
Refinement nodes & $1959$ &
  $1029$ splits, $236$ restricting location updates, $694$ exclusions ($75$ by real location updates, $422$ by
  lower bounds for $\lambda_2$ or $\lambda'$, $197$ by positivity); $4453+1029=4788+694$ &
  \cref{subsec:refinement}, \cref{tab:nodes} &
  Replay; some location rows by separate programs\\
Leaves & $4788$ &
  $3146$ inside ($233$ supplementary) and $1642$ outside; $343$ inside and $70$ outside roots have none &
  \cref{subsec:refinement} &
  Replay; Lean numeric checker\\
Certificate trees & $5591$ &
  One per leaf or subcase of a leaf: $3949$ inside ($349$ for supplementary leaves), $1642$ outside &
  \cref{subsec:subdivision} &
  Replay; independent checker; Lean checker\\
Threshold boxes & $4{,}196{,}879$ &
  Terminal boxes of the trees: $4{,}109{,}455$ inside ($1{,}042{,}363$ for supplementary leaves), $87{,}424$
  outside &
  \cref{def:certificate} &
  As above\\
Relaxation cases & $29{,}397{,}336$ &
  $2^k$ in a box with $k$ tangent rows; $9430$ closed by an exclusion, the rest by integer duals &
  \cref{def:certificate} &
  As above, in exact arithmetic\\
\midrule
\multicolumn{5}{@{}l}{\emph{Rows and exclusions}}\\
\midrule
Near rows & $16{,}842$ &
  All family, shifted and graded rows ($13{,}528$ inside, $3314$ outside), with $11{,}135{,}323$
  entries &
  \cref{subsec:rowkinds} &
  Interval arithmetic; Lean row checker\\
Two-test rows & $37$ &
  On $37$ inside roots: $31$ single, $2$ mixture, $4$ paired &
  \cref{subsec:twotest} &
  Replay; exact comparison\\
Complex location rows & $94$ &
  $25$ for $\lambda_2$ and $69$ for $\lambda'$ (\cref{prop:complexrows,prop:polyrows}) &
  Appendix~\ref{app:rows} &
  Interval arithmetic; separate program\\
Degree-$5$ real rows & $22$ &
  Cells covering $[0.700,0.755]$, used in $300$ location nodes &
  \cref{prop:realrows}, \cref{tab:realrows} &
  Interval arithmetic; replay\\
Positivity exclusions & $197$ &
  Reserved nodes of type $\mathrm{rr}$; least normalized margin $7.8\cdot10^{-5}$ &
  \cref{prop:positivity} &
  Interval arithmetic; replay\\
\bottomrule
\end{tabular}

\medskip
\begin{minipage}{\linewidth}
\footnotesize
Triples of numbers are counts for the types $\mathrm{rr}$, $\mathrm{rc}$, $\mathrm{complex}$;
``supplementary'' refers to the supplementary leaves of \cref{subsec:refinement}. Each split adds one
terminal node; hence the identity in the refinement row. The printed tables were compared with
the pages of \cite{HeathBrown1992zero,Xylouris2011nullstellen}, and Heath-Brown's Tables~4 and~7
were recomputed in floating point (\cref{subsec:locconventions}). The replay does not regenerate
the $94$ complex rows or the conditions \cite[(4.29), (4.34)]{Xylouris2011nullstellen}; separate
interval programs check them. The Lean row checker covers all near rows except the $37$
two-test rows.
\end{minipage}
\caption{The finite analysis of the middle range $0.1\le\lambda_1<1.5$ at a glance. The case
tree, from the parent rows to the leaves, partitions configurations of zeros; the boxes and
relaxation cases partition the thresholds of the near rows. The checks are described in
\cref{sec:verification}.}
\label{tab:census}
\end{table}

\subsection{The leaf model}\label{subsec:leafmodel}
This subsection gives the real coefficients before integer scaling. An interval
$[\lo,\hi)$ is charged at its left endpoint in the objective, where the cost is
largest, and at its right endpoint in density constraints, where the available
weight is smallest. These choices enlarge the feasible set and give an upper bound
for every actual distribution of characters in the interval.

Let $\mathfrak s$ be the specification of a leaf, after the location updates. Let $s$ denote the
\emph{shift} $\min(1.9,p^*,l_2)$, where $p^*=\lo'$ if $\mathfrak s$ has a gap and $p^*=p$ otherwise, and let
$\lambda_3^{\lo}$ be the lower bound for $\lambda_3$ of \cref{prop:third}, computed from $[a,b]$, the
type and, for a reserved $\mathfrak s$, the cap $\hi_2$. The leaf program (\cref{def:leafprogram}) has the
columns of \cref{tab:leafmodel}. Every number is rounded in the safe direction at scale $S$: far weights,
features and the hidden-tail count coefficient $1/w(\mathrm{end})$ down; objective coefficients,
diagonals, correlation terms, the far budget and the constant $\mathrm{first}$ up. Let $O$ be the set
of ordinary characters with a zero in $R_P$, and $\mathcal H$ the set of first-family characters with a
zero in $R_P$.
\begin{enumerate}[(a)]
  \item \emph{Ordinary bins.} Let $R=\max(3,r)$ and let $r=\xi_0<\xi_1<\dots<\xi_N=R$ consist of $r$,
    $R$ and the multiples of $1/\mathrm{den}$ between them, where $\mathrm{den}\in\{200,400,500,800,2000\}$
    is fixed per leaf. The bins are the intervals $[\xi_i,\xi_{i+1})$, except that in a reserved leaf the
    bins with $\xi_{i+1}\le\lambda_3^{\lo}$ are omitted. The ordinary characters beyond $R$ form the tail.
  \item \emph{Reserved family.} Either (fixed terms) the family is charged $n_2\mathcal G_{\phi_2}(\lo_2)$ in
    $\mathrm{first}$ and $n_2w(\hi_2)$ in the far budget, and it appears in the near rows as a family
    term with its feature at $\hi_2$ (there are then no second-family columns, and the count $n_E$ of the
    second-family equality of \cref{def:leafprogram} is $0$); or (columns) the interval $[\lo_2,\hi_2]$ is cut
    at the multiples of $1/400$ into reserved columns $[\lo_{2,j},\hi_{2,j}]$, with values $z_j$ and
    $\sum_jz_j=n_2=n_E$.
  \item \emph{Hidden columns} (outside leaves only). Put $p^\circ=\max(a,p,\lo')$ and
    $\mathrm{end}=\max(3,p^\circ)$. The hidden bins $[\lo_h,\hi_h)$ divide $[p^\circ,\mathrm{end}]$ on the grid
    $1/\mathrm{den}$, and the hidden tail collects the values $t_\chi\ge\mathrm{end}$. The hidden count is at
    most $n$.
\end{enumerate}

\begin{table}[ht]
\centering
\footnotesize
\setlength{\tabcolsep}{3.5pt}
\renewcommand{\arraystretch}{1.35}
\begin{tabular}{@{}>{\raggedright\arraybackslash}p{2.25cm}>{\raggedright\arraybackslash}p{3.3cm}cccccl@{}}
\toprule
Column & Value $x_c$ & $G_c$ & $W_c$ & $C_c$ & $N_c$ & $E_c$ & Near rows\\
\midrule
ordinary bin $[\xi_i,\xi_{i+1})$ & number of $\chi\in O$ with $\nu_{T^*}(\chi)$ in the bin & $\mathcal G(\xi_i)$ &
  $w(\xi_{i+1})$ & $0$ or $1$\textsuperscript{a} & $0$ & $0$ & entry (a)\\
tail & $\sum w(\nu_{T^*}(\chi))$ over $\chi\in O$ with $\nu_{T^*}(\chi)\ge R$ & $\mathcal G(R)/w(R)$ & $1$ & $0$ &
  $0$ & $0$ & omitted\\
reserved column\textsuperscript{b} $[\lo_{2,j},\hi_{2,j}]$ & $z_j=n_2$ if $\nu_*$ is in the column, else $0$ &
  $\mathcal G_{\phi_2}(\lo_{2,j})$ & $w(\hi_{2,j})$ & $0$ & $0$ & $1$ & entry (f)\\
hidden bin\textsuperscript{c} $[\lo_h,\hi_h)$ & number of $\chi\in\mathcal H$ with $t_\chi$ in the bin &
  $e^{-A\lo_h}B_{\phi_1}(p^\circ)$ & $w(\hi_h)$ & $0$ & $1$ & $0$ & entry (g)\\
hidden tail\textsuperscript{c} & $\sum w(t_\chi)$ over $\chi\in\mathcal H$ with $t_\chi\ge\mathrm{end}$ &
  $\dfrac{e^{-A\,\mathrm{end}}B_{\phi_1}(p^\circ)}{w(\mathrm{end})}$ & $1$ & $0$ & $\dfrac1{w(\mathrm{end})}$ & $0$ &
  omitted\\
\bottomrule
\end{tabular}
\caption{The columns of a leaf program and their real coefficients before scaling by $S$:
objective $G_c$, far weight $W_c$, count $C_c$, hidden count $N_c$ and second-family indicator
$E_c$ (\cref{def:leafprogram}). The last column gives the entry of \cref{tab:entries} that supplies
the feature and diagonal of the column in each near row. (a) $1$ if $\mathfrak s$ is unreserved
and $\xi_{i+1}\le\lambda_3^{\lo}$, and $0$ otherwise. (b) Only if the reserved family is charged
through columns. (c) Outside leaves only.}
\label{tab:leafmodel}
\end{table}
Before scaling and rounding, the far budget is $(1+\eta)V-\1_{\mathrm{inside}}\,n\,w(b)$,
with a further subtraction of $n_2w(\hi_2)$ for a fixed reserved-family term.
The stored integer budget $F$ is formed by rounding the initial budget up and each
subtraction down, as in \cref{sec:far}. The count budget is $2$, and the stored
allowance satisfies $\mathrm{final}/S=5\eta$.
For an inside leaf, $\mathrm{first}/S$ is rounded up from
$J_1=\min(J_{\mathrm{old}},J_{\mathrm{new}})$ with anchor $p^*$, omitting
$J_{\mathrm{new}}$ when $p^*<b$. For an outside leaf it starts at $0$.
The fixed reserved-family charge, if used, is added with upward rounding.
The near rows are those of \cref{sec:rows}.

\subsection{Subdivision of a leaf}\label{subsec:subdivision}
A leaf may be certified through finitely many subcases, each of which is certified with its own
leaf model. Besides the second-family split above, three kinds are used. The first replaces the
specification by three specifications. The other two restrict one further quantity $Q$ of the
configuration to one of finitely many intervals $I$ whose union is the range of $Q$; the
configuration set of the subcase $(\mathfrak s,I)$ is the set $C(\mathfrak s,I)$ of configurations for which, for
some admissible choice of $(\chi_1,\rho_1)$, of the minimizers and of $\rho'$, (C1)--(C5) hold and
$Q\in I$. Then $C(\mathfrak s)$ is the union of the sets $C(\mathfrak s,I)$, and all the results of
\cref{sec:rows,sec:casetree} that are stated for $C(\mathfrak s)$ hold for $C(\mathfrak s,I)$ with the same proofs.

\begin{lemma}[Reservation split]\label{lem:reservationsplit}
Let $\mathfrak s$ be unreserved with ordinary lower bound $r$, and let $r<h\le2$. Let $\mathfrak s_1$ and
$\mathfrak s_2$ be $\mathfrak s$ with the reservation $(h,1)$ and $(h,2)$ (so $\lo_2=r$ and $\hi_2=h$), and let $\mathfrak s_0$ be $\mathfrak s$ with
$r$ replaced by $h$. Then $C(\mathfrak s)\subseteq C(\mathfrak s_0)\cup C(\mathfrak s_1)\cup C(\mathfrak s_2)$.
\end{lemma}

\begin{proof}
If $\nu_*<h$, a family $\mathcal F\ne\mathcal F_1$ attaining $\nu_*$ exists, and it is a real character or a pair of
conjugate nonreal characters; the configuration lies in $C(\mathfrak s_1)$ or $C(\mathfrak s_2)$, since
$\nu_*\ge r$. Otherwise every family $\mathcal F\ne\mathcal F_1$ has $\nu(\mathcal F)\ge h$, and the configuration lies in
$C(\mathfrak s_0)$.
\end{proof}

In a reserved child the model uses only the minimality of the reserved family
(\cref{lem:familyfacts}(ii)) and the cap $\lambda_2\le h$ (\cref{lem:familyfacts}(iii)); it never uses
that the reserved family is $\mathcal F_2$. The split of \cref{lem:reservationsplit} is used with $h=\min(\lambda_3^{\lo},2)$ on four inside
roots and on five roots that contain supplementary leaves.

\emph{Height split for type $\mathrm{rc}$.} An inside leaf of type $\mathrm{rc}$ is split by $Q=\mu_1$
(recall that $\mu_1>0$ for this type) into $\mu_1\in[0,1]$ and $\mu_1\in[1,\infty)$, which changes only the
family row (\cref{sec:rows}).

\emph{Second-zero height split.} An inside leaf of type $\mathrm{complex}$ with a finite gap is
split by $Q=|y|$, where $y=\mu_1-\mu_{\rho'}=(\gamma_1-\gamma')\LL$ is the normalized height difference between
$\rho_1$ and the non-distinguished zero $\rho'$ of $\chi_1$ realizing $\lambda'$ (a zero of $\chi_1$, not of $\bchi_1$;
\cref{sec:notation}), into six pieces, according as $|y|$ lies in $[0,1]$, $[1,\frac32]$, $[\frac32,2]$, $[2,3]$,
$[3,5]$ or $[5,\infty)$. This changes only the family row. A piece may be certified by an exclusion (\cref{def:certificate}(i)) rather than by
duals: the near row alone is then infeasible.

\subsection{Realization and the cover theorem}\label{subsec:realization}

\begin{proposition}[Realization]\label{prop:realization}
Let $\mathfrak s$ specify a leaf or one of its subcases, with the program constructed
in \cref{subsec:leafmodel}. For every sufficiently large $q$ whose configuration lies
in $C(\mathfrak s)$, there is a feasible pair $(x,\tau)$ such that
\[
  W(q)+2\eta\le\mathcal V(x).
\]
The column values $x$ are obtained from the actual character counts and weighted tail
masses. The threshold $q_0$ may be chosen uniformly over the finite set of leaves and
subcases.
\end{proposition}

\begin{proof}
Let $O$ and $\mathcal H$ be as in \cref{subsec:leafmodel} ($\mathcal H$ is used only if the
configuration is outside). For
$\chi\in O$ we have $\nu_{T^*}(\chi)\le C_P$ and $\nu_{T^*}(\chi)\ge\nu(\chi)\ge r$ by (C4). Put $x_i$
equal to the number of $\chi\in O$ with $\nu_{T^*}(\chi)\in[\xi_i,\xi_{i+1})$, $x_{\mathrm{tail}}=\sum w(\nu_{T^*}(\chi))$
over $\chi\in O$ with $\nu_{T^*}(\chi)\ge R$, $z_j$ equal to $n_2$ for the reserved column containing
$\nu_*$ and $0$ otherwise, $x^{\mathrm h}_i$ equal to the number of $\chi\in\mathcal H$ whose value $t_\chi$
(\cref{subsec:firstoutside}) lies in the hidden bin $i$, and $x^{\mathrm h}_{\mathrm{tail}}=\sum w(t_\chi)$ over the
$\chi\in\mathcal H$ in the hidden tail. In a reserved leaf no $\chi\in O$ lies in an omitted bin, by \cref{lem:familyfacts}(ii) and
$\lambda_3\ge\lambda_3^{\lo}$.

\emph{Far.} The characters of $O$, the reserved characters, the first-family characters of $\mathcal
H$ and, for an inside configuration, $\chi_1$ and $\bchi_1$ are distinct, and to each we assign one
zero with $|\gamma|\le1$ and parameter at most $C_P$: the $T^*$-representative, the height-one
representative, the zero realizing $t_\chi$, and $\rho_1$ or $\overline{\rho_1}$, which has
$|\gamma_1|\le C_B/\LL\le1$. By \eqref{eq:farbudget}, and since $w$ is decreasing, the far
constraint holds.

\emph{Count.} By \cref{lem:familyfacts}(i), since a character in a counted bin has
$\nu_{T^*}(\chi)<\xi_{i+1}\le\lambda_3^{\lo}\le\lambda_3$.

\emph{Hidden count and second family.} We have $|\mathcal H|\le n$, a character in the hidden
tail contributes $w(t_\chi)/w(\mathrm{end})\le1$ to the hidden count, and $\sum_jz_j=n_2=n_E$ when the reserved
family is charged through columns (otherwise there are no second-family columns and $n_E=0$).

\emph{Near rows.} By \cref{prop:rowsvalid}, each near row holds for some threshold, with features at
least and diagonals at most those of the columns.

\emph{Value.} We use \eqref{eq:Wsplit}. An ordinary character in a bin costs at most the
objective $\mathcal G(\xi_i)$ at the left end of the bin, a reserved character in a column at most $\mathcal G_{\phi_2}$ at
the left end (with fixed terms, at most $\mathcal G_{\phi_2}(\lo_2)$), and the tail characters at most $\mathcal G(R)/w(R)$ per unit of far mass, by the monotonicity
of $\mathcal G$, $\mathcal G_{\phi_2}$ and $\mathcal G/w$ (\cref{lem:envelope}(v)). For an outside configuration we apply
\cref{prop:firstoutside} with the anchor $p^\circ=\max(a,p,\lo')$, which is admissible because every
non-distinguished occurrence of the first family has parameter at least $\lambda'\ge p^\circ$; then
$t_\chi\ge p^\circ$, a character in a hidden bin $[\lo_h,\hi_h)$ costs at most the objective $e^{-A\lo_h}B_{\phi_1}(p^\circ)$
since $e^{-A\lambda}$ is decreasing, and a character in the hidden tail costs at most
$e^{-A\cdot\mathrm{end}}B_{\phi_1}(p^\circ)/w(\mathrm{end})$ per unit of far mass, since $e^{-A\lambda}/w(\lambda)$ is decreasing
(\cref{lem:decay}). Including the directions of integer rounding, we obtain
\[
  W\le\mathcal V(x)-\frac{\mathrm{final}}S+E
    \le\mathcal V(x)-5\eta+2\eta
    \le\mathcal V(x)-2\eta.
\]
\end{proof}

\begin{theorem}[The zero-sum bound in the middle range]\label{thm:cover}
For the case tree and certificates described in this paper, there is a threshold
$q_0$ such that every $q\ge q_0$ with $0.1\le\lambda_1<1.5$ satisfies
\[
  W(q)<1-2\eta.
\]
More precisely, its configuration belongs to a leaf or subcase whose program has a
feasible pair $(x,\tau)$ with $W(q)+2\eta\le\mathcal V(x)<1$.
\end{theorem}

\begin{proof}
By \cref{thm:sourcecover} the configuration lies in the configuration set of a root. Following the
refinement tree and the subdivisions of \cref{subsec:subdivision}, each of which is exhaustive,
we reach a leaf or subcase whose configuration set contains the configuration; it is not an
excluded node, because excluded nodes have no configurations. \Cref{prop:realization} gives the
feasible pair, and \cref{thm:certificate} and the certificates of \cref{sec:verification} give
$\mathcal V(x)<1$.
\end{proof}
